\section{A Causal Policy That Retains Observations}\label{sec:policy}
\subsection{Clock Changes and Lifetime Statistics}
For each available arm, store its lifetime sample count $N_i$ and reward sum $S_i$. Set an integer variable $\mathrm{base}$ to zero initially. If only one arm is available, pull it and update its statistics. Otherwise, after receiving the current batch, start a new clock segment whenever
\[
 \mathrm{base}=0\quad\text{or}\quad K_t-1\ge2\,\mathrm{base}.
\]
At that time set
\begin{equation}\label{eq:clock}
 \mathrm{base}=K_t-1,\qquad \sigma=t,\qquad d=2\,\mathrm{base}.
\end{equation}
At all other times retain $\sigma$ and $d$. The local clock is $s=t-\sigma+1$. Select an arm maximizing
\begin{equation}\label{eq:index}
 U_i(t)=
 \begin{cases}
 +\infty,&N_i(t-1)=0,\\[1mm]
 \widehat\mu_{i,N_i(t-1)}+
 \displaystyle\sqrt{\frac{[\log\{s/(dN_i(t-1))\}]_+}{N_i(t-1)}},
 &N_i(t-1)\ge1,
 \end{cases}
\end{equation}
where $\widehat\mu_{i,n}$ is the average of the first $n$ rewards from arm $i$. Break ties by arrival order and a fixed ordering within a batch. Observe one reward, then increment the chosen arm's lifetime count and sum. No clock change modifies either statistic. We call this policy MC-MOSS, for MOSS with a recalibrated clock and retained memory.

Every decision is well defined even if a batch is too large to initialize fully. A batch that crosses a threshold belongs to the new segment, since \eqref{eq:clock} precedes the action. If the segment begins with $c$ competing arms, its total number of available arms stays at most $d=2c$: otherwise a new segment would have begun before the next action. This capacity bound holds without predicting the final arm count.

\subsection{Why Old Observations Can Be Used}
The policy does not restart a statistical experiment at a clock change. Conditional on the observed past, old empirical means need not have the distribution of fresh sample averages. We therefore work on the original potential reward streams and take unconditional expectations throughout the local argument.

For one stream of mean $\mu$, let
\begin{equation}\label{eq:underestimate}
 Z(v)=\sup_{n\ge1}\left[\mu-\widehat\mu_n-
 \sqrt{[\log(v/n)]_+/n}\right]_+.
\end{equation}
A maximal concentration argument gives
\begin{equation}\label{eq:Zbounds}
 \EE Z(v)\le5v^{-1/2},\qquad
 \PP\{Z(v)>x\}\le\min\{1,5/(vx^2)\}.
\end{equation}
At each physical round, choose a deterministic label for a best available arm, using the fixed means and schedule. On every action selecting an observed arm, index comparison yields
\begin{equation}\label{eq:index-decomposition}
 \Delta_{I_t}(t)\le Z_{i_t^*}(s/d)
 +(\widehat\mu_{I_t,n}-\mu_{I_t})_+
 +\sqrt{[\log\{s/(dn)\}]_+/n},\qquad n=N_{I_t}(t-1).
\end{equation}
If the best arm is still unobserved, the infinite-index rule selects an unobserved arm, which is charged separately.

For the other terms, each pair consisting of an arm and its count just before a pull occurs at most once. If at most $h$ such pulls use $k$ arms, this pathwise fact gives
\[
 \EE\sum(\widehat\mu_{I_t,n}-\mu_{I_t})_+\le\sqrt{2kh},
 \qquad
 \sum_{n\ge1}\sqrt{[\log(v/n)]_+/n}\le\sqrt{2\pi v}.
\]
No independence between a random count and its empirical mean is used. Old samples may cause some counts to be skipped, which can only reduce these upper sums.

\begin{lemma}[A segment with retained history]\label{lem:kernel}
Consider a segment with deterministic starting time and capacity $d$, using \eqref{eq:index} and all rewards gathered in the experiment. For any deterministic prefix of length $h$ within the segment, its expected regret satisfies
\begin{equation}\label{eq:kernel}
 \EE R_{\mathrm{segment}}\le\min\{h,15\sqrt{dh}\}.
\end{equation}
Arrivals within the segment are allowed, provided at most $d$ arms are available. The expectation is under the original experiment, not conditional on an arbitrary realized history.
\end{lemma}

The terms above sum to at most $(1+10+\sqrt2+\sqrt{2\pi})\sqrt{dh}$. The extra one accounts for genuine first observations. To pass from segments to \eqref{eq:main}, note that their initial competing-arm counts grow geometrically. Group the segment starts by remaining-lifetime scale, and assign each segment to the arms that caused its capacity increase. Delaying these arms to the segment start changes the lifetime spectrum by at most a factor of two. This charges the sum of the local square-root costs to $\cS_T$. Once the competing-arm count exceeds the remaining horizon, the entire tail is charged by its length. Full details and the lower-bound construction appear in the supplement.

\subsection{Instance Behavior and the Role of Memory}
For each arm that is strictly suboptimal at some time before $T$, define
\[
 \gamma_i=\min_{t\le T:\,i\in A_t,\,\Delta_i(t)>0}\Delta_i(t),
 \qquad \gamma=\min_i\gamma_i.
\]
These positive quantities are defined retrospectively and are not inputs to the policy. They need not be bounded away from zero as the experiment grows.

\begin{proposition}[Gap bounds and a paired restart comparison]\label{prop:gap}
If at least one positive gap occurs, MC-MOSS satisfies
\begin{equation}\label{eq:gap-bound}
 \EE R_T\le\min\left\{T,300\cS_T,
 \frac{400(K_T-1)\log(eT)}{\gamma}\right\}.
\end{equation}
For deterministic rewards, the stronger pathwise bound
\begin{equation}\label{eq:det-gap}
 R_T\le\sum_{i:\gamma_i\text{ defined}}
 \left\{1+\frac{2\log(eT)}{\gamma_i}\right\}
\end{equation}
holds. There is a deterministic fixed-gap schedule family on which
\[
 R_T(\pi^{\mathrm{MC}})=O(\log T),\qquad
 R_T(\pi^{\mathrm{cold}})=\Omega\{(\log T)^2\},
\]
where $\pi^{\mathrm{cold}}$ has the same indices, clock changes and tie rules, but erases all statistics at each clock change. Both policies obey the upper rate in Theorem~\ref{thm:main}.
\end{proposition}

The separating family has one inferior old arm and newly arriving arms tied with the incumbent. Erasing observations repeatedly removes the evidence that the old arm is inferior. The proof counts actual inferior-arm pulls; it does not compare two regret upper bounds. This comparison isolates a cost of forgetting. It makes no assertion of dominance over AUER, ordinary all-history UCB, or other MOSS variants.

A simple boundary example further separates minimax and instance behavior. Suppose the initial arm always gives reward one, and one arm giving reward zero arrives at every later round. The infinite-index rule explores every new arm and incurs regret $T-1$. Capping observed indices at one and assigning index one to unobserved arms, with the same tie order, makes regret zero after the initial observation. The capped policy retains the bounds in Theorem~\ref{thm:main} and Proposition~\ref{prop:gap}; its proof is given in the supplement. The worst-case risk for this arrival schedule is linear, so both behaviors are compatible with minimax optimality.
